\documentclass[11pt,a4paper]{article}

\usepackage[a4paper,margin=20mm]{geometry}
\usepackage{amsmath,amssymb}
\usepackage{booktabs,array}
\usepackage[hidelinks]{hyperref}

\setlength{\parskip}{6pt plus 2pt minus 1pt}
\setlength{\parindent}{0pt}

\newcommand{\optsfourA}[4]{%
  \begin{tabbing}
    \hspace{0.48\linewidth}\=\kill
    (A) #1 \> (B) #2 \\[2pt]
    (C) #3 \> (D) #4
  \end{tabbing}%
}

\newcommand{\optslistA}[4]{%
  \begin{itemize}
    \setlength{\itemsep}{2pt}\setlength{\parskip}{0pt}
    \item[(A)] #1
    \item[(B)] #2
    \item[(C)] #3
    \item[(D)] #4
  \end{itemize}%
}

\newcommand{\ans}[1]{\textbf{Answer:} #1\par\medskip}

\begin{document}

% ============================================================================
% GATE 2023 -- GENERAL APTITUDE (GA) & MATHEMATICS (MA)
% ============================================================================

\begin{center}
  {\LARGE\bfseries GATE 2023}\\[4pt]
  {\large\bfseries General Aptitude (GA) \& Mathematics (MA)}\\[2pt]
  {\small\itshape Official Final Answer Key attached to every question}
\end{center}
\medskip

\section*{General Aptitude (GA)}

\subsection*{Q.1 -- Q.5 Multiple Choice Question (MCQ), carry ONE mark each (for each wrong answer: $-1/3$).}

\textbf{Q.1} The village was nestled in a green spot, \underline{\hspace{2cm}} the ocean and the hills.
\optsfourA{through}{in}{at}{between}
\ans{(D)}

\textbf{Q.2} Disagree : Protest : : Agree : \underline{\hspace{2cm}} (By word meaning)
\optsfourA{Refuse}{Pretext}{Recommend}{Refute}
\ans{(C)}

\textbf{Q.3} A `frabjous' number is defined as a 3 digit number with all digits odd, and no two adjacent digits being the same. For example, $137$ is a frabjous number, while $133$ is not. How many such frabjous numbers exist?
\optsfourA{$125$}{$720$}{$60$}{$80$}
\ans{(D)}

\textbf{Q.4} Which one among the following statements must be TRUE about the mean and the median of the scores of all candidates appearing for GATE 2023?
\optslistA{The median is at least as large as the mean.}{The mean is at least as large as the median.}{At most half the candidates have a score that is larger than the median.}{At most half the candidates have a score that is larger than the mean.}
\ans{(C)}

\textbf{Q.5} In the given diagram, ovals are marked at different heights ($h$) of a hill. Which one of the following options P, Q, R, and S depicts the top view of the hill?\\[4pt]
\emph{[Diagram note: Side profile of an asymmetric hill peaking at distance $\approx 0.22\text{ km}$ (height $\approx 0.78\text{ km}$), with a steep slope on the left ($0\text{ km}$ side) and a gentler, longer slope on the right (extending past $0.8\text{ km}$), with horizontal oval cross-sections at several altitudes. In the four top-view contour diagrams with $0\text{ km}$ on the left and distance increasing rightward: P has closely spaced contours on the right (peak shifted right); Q has closely spaced contours on the left near $0\text{ km}$ and wider spacing toward the right (peak shifted left); R has concentric symmetric contours; S has non-nested/offset contours shifted right.]}
\optsfourA{P}{Q}{R}{S}
\ans{(B)}

\subsection*{Q.6 -- Q.10 Multiple Choice Question (MCQ), carry TWO marks each (for each wrong answer: $-2/3$).}

\textbf{Q.6} Residency is a famous housing complex with many well-established individuals among its residents. A recent survey conducted among the residents of the complex revealed that all of those residents who are well established in their respective fields happen to be academicians. The survey also revealed that most of these academicians are authors of some best-selling books.\\
Based only on the information provided above, which one of the following statements can be logically inferred with certainty?
\optslistA{Some residents of the complex who are well established in their fields are also authors of some best-selling books.}{All academicians residing in the complex are well established in their fields.}{Some authors of best-selling books are residents of the complex who are well established in their fields.}{Some academicians residing in the complex are well established in their fields.}
\ans{MTA (Marks to All)}

\textbf{Q.7} Ankita has to climb $5$ stairs starting at the ground, while respecting the following rules:
\begin{enumerate}
  \setlength{\itemsep}{1pt}\setlength{\parskip}{0pt}
  \item At any stage, Ankita can move either one or two stairs up.
  \item At any stage, Ankita cannot move to a lower step.
\end{enumerate}
Let $F(N)$ denote the number of possible ways in which Ankita can reach the $N$th stair. For example, $F(1) = 1$, $F(2) = 2$, $F(3) = 3$. The value of $F(5)$ is \underline{\hspace{2cm}}.
\optsfourA{$8$}{$7$}{$6$}{$5$}
\ans{(A)}

\textbf{Q.8} The information contained in DNA is used to synthesize proteins that are necessary for the functioning of life. DNA is composed of four nucleotides: Adenine (A), Thymine (T), Cytosine (C), and Guanine (G). The information contained in DNA can then be thought of as a sequence of these four nucleotides: A, T, C, and G. DNA has coding and non-coding regions. Coding regions---where the sequence of these nucleotides are read in groups of three to produce individual amino acids---constitute only about $2\%$ of human DNA. For example, the triplet of nucleotides CCG codes for the amino acid glycine, while the triplet GGA codes for the amino acid proline. Multiple amino acids are then assembled to form a protein.\\
Based only on the information provided above, which of the following statements can be logically inferred with certainty?
\begin{itemize}
  \setlength{\itemsep}{1pt}\setlength{\parskip}{0pt}
  \item[(i)] The majority of human DNA has no role in the synthesis of proteins.
  \item[(ii)] The function of about $98\%$ of human DNA is not understood.
\end{itemize}
\optsfourA{only (i)}{only (ii)}{both (i) and (ii)}{neither (i) nor (ii)}
\ans{(D)}

\textbf{Q.9} Which one of the given figures P, Q, R and S represents the graph of the following function?
\[
  f(x) = \bigl|\,|x + 2| - |x - 1|\,\bigr|
\]
\emph{[Diagram note: Four piecewise-linear graphs are shown: P is constant at $3$ for $x \le -2$, decreases linearly to $0$ at $x = -0.5$, increases linearly to $3$ at $x = 1$, and stays constant at $3$ for $x \ge 1$; Q is $-3$ for $x \le -2$, rises linearly through $(-0.5, 0)$ to $3$ at $x = 1$, and stays at $3$; R rises linearly to the left of $x = -0.5$ (reaching $9$ at $x = -5$), hits $0$ at $x = -0.5$, rises to $3$ at $x = 1$, and stays at $3$; S descends linearly to $3$ at $x = -2$ and stays constant at $3$ for $x \ge -2$.]}
\optsfourA{P}{Q}{R}{S}
\ans{(A)}

\textbf{Q.10} An opaque cylinder is suspended in the path of a parallel beam of light, such that its shadow is cast on a screen oriented perpendicular to the direction of the light beam. The cylinder can be reoriented in any direction within the light beam. Under these conditions, which one of the shadows P, Q, R, and S is NOT possible?\\[4pt]
\emph{[Diagram note: Four grey shadow shapes are shown: P is a solid circle; Q is a stadium/capsule shape (rectangle with two semicircular ends); R is a rectangle; S is a non-rectangular parallelogram with sharp acute and obtuse corners.]}
\optsfourA{P}{Q}{R}{S}
\ans{(D)}

\section*{Mathematics (MA)}

\textbf{Useful Data:} $\mathbb{N}$: set of all positive integers; $\mathbb{Z}$: set of all integers; $\mathbb{Q}$: set of all rational numbers; $\mathbb{R}$: set of all real numbers; $\mathbb{C}$: set of all complex numbers; $\mathbb{R}^n = \{(x_1, x_2, \ldots, x_n) : x_i \in \mathbb{R}, \, i = 1, 2, \ldots, n\}$; $\mathbb{R} \times \mathbb{R}$: Cartesian product of $\mathbb{R}$ with $\mathbb{R}$.

\subsection*{Q.11 -- Q.24 Multiple Choice Question (MCQ), carry ONE mark each (for each wrong answer: $-1/3$).}

\textbf{Q.11} Let $f, g : \mathbb{R}^2 \to \mathbb{R}$ be defined by
\[
  f(x, y) = x^2 - \frac{3}{2} x y^2 \quad \text{and} \quad g(x, y) = 4x^4 - 5x^2 y + y^2
\]
for all $(x, y) \in \mathbb{R}^2$. Consider the following statements:\\
$P$: $f$ has a saddle point at $(0, 0)$.\\
$Q$: $g$ has a saddle point at $(0, 0)$.\\
Then
\optslistA{both $P$ and $Q$ are TRUE}{$P$ is FALSE but $Q$ is TRUE}{$P$ is TRUE but $Q$ is FALSE}{both $P$ and $Q$ are FALSE}
\ans{(A)}

\textbf{Q.12} Let $\mathbb{R}^3$ be a topological space with the usual topology and $\mathbb{Q}$ denote the set of rational numbers. Define the subspaces $X, Y, Z$ and $W$ of $\mathbb{R}^3$ as follows:
\begin{align*}
  X &= \{(x, y, z) \in \mathbb{R}^3 : |x| + |y| + |z| \in \mathbb{Q}\}, \\
  Y &= \{(x, y, z) \in \mathbb{R}^3 : xyz = 1\}, \\
  Z &= \{(x, y, z) \in \mathbb{R}^3 : x^2 + y^2 + z^2 = 1\}, \\
  W &= \{(x, y, z) \in \mathbb{R}^3 : xyz = 0\}.
\end{align*}
Which of the following statements is correct?
\optsfourA{$X$ is homeomorphic to $Y$}{$Z$ is homeomorphic to $W$}{$Y$ is homeomorphic to $W$}{$X$ is NOT homeomorphic to $W$}
\ans{(D)}

\textbf{Q.13} Let $P(x) = 1 + e^{2\pi i x} + 2 e^{3\pi i x}$, $x \in \mathbb{R}$, $i = \sqrt{-1}$. Then
\[
  \lim_{N \to \infty} \frac{1}{N} \sum_{k=0}^{N-1} P(k\sqrt{2})
\]
is equal to
\optsfourA{$0$}{$1$}{$3$}{$4$}
\ans{(B)}

\textbf{Q.14} Let $T : \mathbb{R}^3 \to \mathbb{R}^3$ be a linear transformation satisfying
\[
  T(1, 0, 0) = (0, 1, 1), \quad T(1, 1, 0) = (1, 0, 1) \quad \text{and} \quad T(1, 1, 1) = (1, 1, 2).
\]
Then
\optslistA{$T$ is one-one but $T$ is NOT onto}{$T$ is one-one and onto}{$T$ is NEITHER one-one NOR onto}{$T$ is NOT one-one but $T$ is onto}
\ans{(C)}

\textbf{Q.15} Let $\mathbb{D} = \{z \in \mathbb{C} : |z| < 1\}$ and $f : \mathbb{D} \to \mathbb{C}$ be defined by
\[
  f(z) = z - 25z^3 + \frac{z^5}{5!} - \frac{z^7}{7!} + \frac{z^9}{9!} - \frac{z^{11}}{11!}.
\]
Consider the following statements:\\
$P$: $f$ has three zeros (counting multiplicity) in $\mathbb{D}$.\\
$Q$: $f$ has one zero in $\mathbb{U} = \{z \in \mathbb{C} : \frac{1}{2} < |z| < 1\}$.\\
Then
\optslistA{$P$ is TRUE but $Q$ is FALSE}{$P$ is FALSE but $Q$ is TRUE}{both $P$ and $Q$ are TRUE}{both $P$ and $Q$ are FALSE}
\ans{(A)}

\textbf{Q.16} Let $\mathcal{N} \subseteq \mathbb{R}$ be a non-measurable set with respect to the Lebesgue measure on $\mathbb{R}$. Consider the following statements:\\
$P$: If $M = \{x \in \mathcal{N} : x \text{ is irrational}\}$, then $M$ is Lebesgue measurable.\\
$Q$: The boundary of $\mathcal{N}$ has positive Lebesgue outer measure.\\
Then
\optslistA{both $P$ and $Q$ are TRUE}{$P$ is FALSE and $Q$ is TRUE}{$P$ is TRUE and $Q$ is FALSE}{both $P$ and $Q$ are FALSE}
\ans{(B)}

\textbf{Q.17} For $k \in \mathbb{N}$, let $E_k$ be a measurable subset of $[0, 1]$ with Lebesgue measure $\frac{1}{k^2}$. Define
\[
  E = \bigcap_{n=1}^{\infty} \bigcup_{k=n}^{\infty} E_k \quad \text{and} \quad F = \bigcup_{n=1}^{\infty} \bigcap_{k=n}^{\infty} E_k.
\]
Consider the following statements:\\
$P$: Lebesgue measure of $E$ is equal to zero.\\
$Q$: Lebesgue measure of $F$ is equal to zero.\\
Then
\optslistA{both $P$ and $Q$ are TRUE}{both $P$ and $Q$ are FALSE}{$P$ is TRUE but $Q$ is FALSE}{$Q$ is TRUE but $P$ is FALSE}
\ans{(A)}

\textbf{Q.18} Consider $\mathbb{R}^2$ with the usual Euclidean metric. Let
\begin{align*}
  X &= \left\{\left(x, x \sin\frac{1}{x}\right) \in \mathbb{R}^2 : x \in (0, 1]\right\} \cup \{(0, y) \in \mathbb{R}^2 : -\infty < y < \infty\} \quad \text{and} \\
  Y &= \left\{\left(x, \sin\frac{1}{x}\right) \in \mathbb{R}^2 : x \in (0, 1]\right\} \cup \{(0, y) \in \mathbb{R}^2 : -\infty < y < \infty\}.
\end{align*}
Consider the following statements:\\
$P$: $X$ is a connected subset of $\mathbb{R}^2$.\\
$Q$: $Y$ is a connected subset of $\mathbb{R}^2$.\\
Then
\optslistA{both $P$ and $Q$ are TRUE}{$P$ is FALSE and $Q$ is TRUE}{$P$ is TRUE and $Q$ is FALSE}{both $P$ and $Q$ are FALSE}
\ans{(A)}

\textbf{Q.19} Let $M = \begin{bmatrix} 4 & -3 \\ 1 & 0 \end{bmatrix}$. Consider the following statements:\\
$P$: $M^8 + M^{12}$ is diagonalizable.\\
$Q$: $M^7 + M^9$ is diagonalizable.\\
Which of the following statements is correct?
\optslistA{$P$ is TRUE and $Q$ is FALSE}{$P$ is FALSE and $Q$ is TRUE}{Both $P$ and $Q$ are FALSE}{Both $P$ and $Q$ are TRUE}
\ans{(D)}

\textbf{Q.20} Let $C[0, 1] = \{f : [0, 1] \to \mathbb{R} : f \text{ is continuous}\}$. Consider the metric space $(C[0, 1], d_\infty)$, where
\[
  d_\infty(f, g) = \sup\{|f(x) - g(x)| : x \in [0, 1]\} \quad \text{for } f, g \in C[0, 1].
\]
Let $f_0(x) = 0$ for all $x \in [0, 1]$ and
\[
  X = \left\{f \in (C[0, 1], d_\infty) : d_\infty(f_0, f) \ge \frac{1}{2}\right\}.
\]
Let $f_1, f_2 \in C[0, 1]$ be defined by $f_1(x) = x$ and $f_2(x) = 1 - x$ for all $x \in [0, 1]$. Consider the following statements:\\
$P$: $f_1$ is in the interior of $X$.\\
$Q$: $f_2$ is in the interior of $X$.\\
Which of the following statements is correct?
\optslistA{$P$ is TRUE and $Q$ is FALSE}{$P$ is FALSE and $Q$ is TRUE}{Both $P$ and $Q$ are FALSE}{Both $P$ and $Q$ are TRUE}
\ans{(D)}

\textbf{Q.21} Consider the metrics $\rho_1$ and $\rho_2$ on $\mathbb{R}$, defined by
\[
  \rho_1(x, y) = |x - y| \quad \text{and} \quad \rho_2(x, y) = \begin{cases} 0, & \text{if } x = y, \\ 1, & \text{if } x \ne y. \end{cases}
\]
Let $X = \{n \in \mathbb{N} : n \ge 3\}$ and $Y = \left\{n + \frac{1}{n} : n \in \mathbb{N}\right\}$. Define $f : X \cup Y \to \mathbb{R}$ by
\[
  f(x) = \begin{cases} 2, & \text{if } x \in X, \\ 3, & \text{if } x \in Y. \end{cases}
\]
Consider the following statements:\\
$P$: The function $f : (X \cup Y, \rho_1) \to (\mathbb{R}, \rho_1)$ is uniformly continuous.\\
$Q$: The function $f : (X \cup Y, \rho_2) \to (\mathbb{R}, \rho_1)$ is uniformly continuous.\\
Then
\optslistA{$P$ is TRUE and $Q$ is FALSE}{$P$ is FALSE and $Q$ is TRUE}{both $P$ and $Q$ are FALSE}{both $P$ and $Q$ are TRUE}
\ans{(B)}

\textbf{Q.22} Let $T : \mathbb{R}^4 \to \mathbb{R}^4$ be a linear transformation and the null space of $T$ be the subspace of $\mathbb{R}^4$ given by
\[
  \{(x_1, x_2, x_3, x_4) \in \mathbb{R}^4 : 4x_1 + 3x_2 + 2x_3 + x_4 = 0\}.
\]
If $\operatorname{Rank}(T - 3I) = 3$, where $I$ is the identity map on $\mathbb{R}^4$, then the minimal polynomial of $T$ is
\optsfourA{$x(x - 3)$}{$x(x - 3)^3$}{$x^3(x - 3)$}{$x^2(x - 3)^2$}
\ans{(A)}

\textbf{Q.23} Let $C[0, 1]$ denote the set of all real valued continuous functions defined on $[0, 1]$ and $\|f\|_\infty = \sup\{|f(x)| : x \in [0, 1]\}$ for all $f \in C[0, 1]$. Let
\[
  X = \{f \in C[0, 1] : f(0) = f(1) = 0\}.
\]
Define $F : (C[0, 1], \|\cdot\|_\infty) \to \mathbb{R}$ by $F(f) = \int_0^1 f(t)\,dt$ for all $f \in C[0, 1]$. Denote $S_X = \{f \in X : \|f\|_\infty = 1\}$. Then the set $\{f \in X : F(f) = \|F\|\} \cap S_X$ has
\optslistA{NO element}{exactly one element}{exactly two elements}{an infinite number of elements}
\ans{(A)}

\textbf{Q.24} Let $X$ and $Y$ be two topological spaces. A continuous map $f : X \to Y$ is said to be proper if $f^{-1}(K)$ is compact in $X$ for every compact subset $K$ of $Y$, where $f^{-1}(K)$ is defined by $f^{-1}(K) = \{x \in X : f(x) \in K\}$. Consider $\mathbb{R}$ with the usual topology. If $\mathbb{R} \setminus \{0\}$ has the subspace topology induced from $\mathbb{R}$ and $\mathbb{R} \times \mathbb{R}$ has the product topology, then which of the following maps is proper?
\optslistA{$f : \mathbb{R} \setminus \{0\} \to \mathbb{R}$ defined by $f(x) = x$}{$f : \mathbb{R} \times \mathbb{R} \to \mathbb{R} \times \mathbb{R}$ defined by $f(x, y) = (x + y, y)$}{$f : \mathbb{R} \times \mathbb{R} \to \mathbb{R}$ defined by $f(x, y) = x$}{$f : \mathbb{R} \times \mathbb{R} \to \mathbb{R}$ defined by $f(x, y) = x^2 - y^2$}
\ans{(B)}

\subsection*{Q.25 -- Q.35 Numerical Answer Type (NAT), carry ONE mark each (no negative marking).}

\textbf{Q.25} Consider the following Linear Programming Problem P:
\begin{align*}
  \text{Minimize} \quad & 3x_1 + 4x_2 \\
  \text{subject to} \quad & x_1 - x_2 \le 1, \\
  & x_1 + x_2 \ge 3, \\
  & x_1 \ge 0, \quad x_2 \ge 0.
\end{align*}
The optimal value of the problem P is \underline{\hspace{2cm}}.
\ans{$10$}

\textbf{Q.26} Let $u(x, t)$ be the solution of
\begin{align*}
  \frac{\partial^2 u}{\partial x^2} - \frac{1}{c^2} \frac{\partial^2 u}{\partial t^2} &= 0, \quad x \in (-\infty, \infty), \quad t > 0, \\
  u(x, 0) &= \sin x, \quad x \in (-\infty, \infty), \\
  \frac{\partial u}{\partial t}(x, 0) &= \cos x, \quad x \in (-\infty, \infty),
\end{align*}
for some positive real number $c$. Let the domain of dependence of the solution $u$ at the point $P(3, 2)$ be the line segment on the $x$-axis with end points $Q$ and $R$. If the area of the triangle $PQR$ is $8$ square units, then the value of $c^2$ is \underline{\hspace{2cm}}.
\ans{$4$}

\textbf{Q.27} Let
\[
  \frac{z}{1 - z - z^2} = \sum_{n=0}^{\infty} a_n z^n, \quad a_n \in \mathbb{R}
\]
for all $z$ in some neighbourhood of $0$ in $\mathbb{C}$. Then the value of $a_6 + a_5$ is equal to \underline{\hspace{2cm}}.
\ans{$13$}

\textbf{Q.28} Let $p(x) = x^3 - 2x + 2$. If $q(x)$ is the interpolating polynomial of degree less than or equal to $4$ for the data
\begin{center}
\begin{tabular}{cccccc}
\toprule
$x$ & $-2$ & $-1$ & $0$ & $1$ & $3$ \\
\midrule
$q(x)$ & $p(-2)$ & $p(-1)$ & $2.5$ & $p(1)$ & $p(3)$ \\
\bottomrule
\end{tabular}
\end{center}
then the value of $\dfrac{d^4 q}{dx^4}$ at $x = 0$ is \underline{\hspace{2cm}}.
\ans{$2$}

\textbf{Q.29} For a fixed $c \in \mathbb{R}$, let $\alpha = \int_0^2 (9x^2 - 5cx^4)\,dx$. If the value of $\int_0^2 (9x^2 - 5cx^4)\,dx$ obtained by using the Trapezoidal rule is equal to $\alpha$, then the value of $c$ is \underline{\hspace{2cm}} (rounded off to 2 decimal places).
\ans{$0.24 \text{ to } 0.26$}

\textbf{Q.30} If for some $\alpha \in \mathbb{R}$,
\[
  \int_1^4 \int_{-x}^{x} \frac{1}{x^2 + y^2}\,dy\,dx = \int_{-\pi/4}^{\pi/4} \int_{\sec\theta}^{\alpha\sec\theta} \frac{1}{r}\,dr\,d\theta,
\]
then the value of $\alpha$ equals \underline{\hspace{2cm}}.
\ans{$4$}

\textbf{Q.31} Let $S$ be the portion of the plane $z = 2x + 2y - 100$ which lies inside the cylinder $x^2 + y^2 = 1$. If the surface area of $S$ is $\alpha\pi$, then the value of $\alpha$ is equal to \underline{\hspace{2cm}}.
\ans{$3$}

\textbf{Q.32} Let
\[
  L^2[-1, 1] = \left\{f : [-1, 1] \to \mathbb{R} : f \text{ is Lebesgue measurable and } \int_{-1}^{1} |f(x)|^2\,dx < \infty\right\}
\]
and the norm $\|f\|_2 = \left(\int_{-1}^{1} |f(x)|^2\,dx\right)^{1/2}$ for $f \in L^2[-1, 1]$. Let $F : (L^2[-1, 1], \|\cdot\|_2) \to \mathbb{R}$ be defined by
\[
  F(f) = \int_{-1}^{1} f(x) x^2\,dx \quad \text{for all } f \in L^2[-1, 1].
\]
If $\|F\|$ denotes the norm of the linear functional $F$, then $5\|F\|^2$ is equal to \underline{\hspace{2cm}}.
\ans{$2$}

\textbf{Q.33} Let $y(t)$ be the solution of the initial value problem
\[
  y'' + 4y = \begin{cases} t, & 0 \le t \le 2, \\ 2, & 2 < t < \infty, \end{cases} \quad \text{and} \quad y(0) = y'(0) = 0.
\]
If $\alpha = y\!\left(\frac{\pi}{2}\right)$, then the value of $\frac{4}{\pi}\alpha$ is \underline{\hspace{2cm}} (rounded off to 2 decimal places).
\ans{$0.49 \text{ to } 0.51$}

\textbf{Q.34} Consider $\mathbb{R}^4$ with the inner product $\langle x, y \rangle = \sum_{i=1}^4 x_i y_i$, for $x = (x_1, x_2, x_3, x_4)$ and $y = (y_1, y_2, y_3, y_4)$. Let $M = \{(x_1, x_2, x_3, x_4) \in \mathbb{R}^4 : x_1 = x_3\}$ and $M^\perp$ denote the orthogonal complement of $M$. The dimension of $M^\perp$ is equal to \underline{\hspace{2cm}}.
\ans{$1$}

\textbf{Q.35} Let
\[
  M = \begin{bmatrix} 3 & -1 & -2 \\ 0 & 2 & 4 \\ 0 & 0 & 1 \end{bmatrix} \quad \text{and} \quad I = \begin{bmatrix} 1 & 0 & 0 \\ 0 & 1 & 0 \\ 0 & 0 & 1 \end{bmatrix}.
\]
If $6M^{-1} = M^2 - 6M + \alpha I$ for some $\alpha \in \mathbb{R}$, then the value of $\alpha$ is equal to \underline{\hspace{2cm}}.
\ans{$11$}

\subsection*{Q.36 -- Q.41 Multiple Choice Question (MCQ), carry TWO marks each (for each wrong answer: $-2/3$).}

\textbf{Q.36} Let $GL_2(\mathbb{C})$ denote the group of $2 \times 2$ invertible complex matrices with usual matrix multiplication. For $S, T \in GL_2(\mathbb{C})$, $\langle S, T \rangle$ denotes the subgroup generated by $S$ and $T$. Let $S = \begin{bmatrix} 0 & -1 \\ 1 & 0 \end{bmatrix} \in GL_2(\mathbb{C})$ and $G_1, G_2, G_3$ be three subgroups of $GL_2(\mathbb{C})$ given by
\begin{align*}
  G_1 &= \langle S, T_1 \rangle, \quad \text{where } T_1 = \begin{bmatrix} i & 0 \\ 0 & i \end{bmatrix}, \\
  G_2 &= \langle S, T_2 \rangle, \quad \text{where } T_2 = \begin{bmatrix} i & 0 \\ 0 & -i \end{bmatrix}, \\
  G_3 &= \langle S, T_3 \rangle, \quad \text{where } T_3 = \begin{bmatrix} 0 & 1 \\ 1 & 0 \end{bmatrix}.
\end{align*}
Let $Z(G_i)$ denote the center of $G_i$ for $i = 1, 2, 3$. Which of the following statements is correct?
\optsfourA{$G_1$ is isomorphic to $G_3$}{$Z(G_1)$ is isomorphic to $Z(G_2)$}{$Z(G_3) = \left\{\begin{bmatrix} 1 & 0 \\ 0 & 1 \end{bmatrix}\right\}$}{$Z(G_2)$ is isomorphic to $Z(G_3)$}
\ans{(D)}

\textbf{Q.37} Let $\ell^2 = \{(x_1, x_2, x_3, \ldots) : x_n \in \mathbb{R} \text{ for all } n \in \mathbb{N} \text{ and } \sum_{n=1}^{\infty} x_n^2 < \infty\}$. For a sequence $(x_1, x_2, x_3, \ldots) \in \ell^2$, define $\|(x_1, x_2, x_3, \ldots)\|_2 = \left(\sum_{n=1}^{\infty} x_n^2\right)^{1/2}$. Let $S : (\ell^2, \|\cdot\|_2) \to (\ell^2, \|\cdot\|_2)$ and $T : (\ell^2, \|\cdot\|_2) \to (\ell^2, \|\cdot\|_2)$ be defined by
\begin{align*}
  S(x_1, x_2, x_3, \ldots) &= (y_1, y_2, y_3, \ldots), \quad \text{where } y_n = \begin{cases} 0, & n = 1, \\ x_{n-1}, & n \ge 2, \end{cases} \\
  T(x_1, x_2, x_3, \ldots) &= (y_1, y_2, y_3, \ldots), \quad \text{where } y_n = \begin{cases} 0, & n \text{ is odd}, \\ x_n, & n \text{ is even}. \end{cases}
\end{align*}
Then
\optslistA{$S$ is a compact linear map and $T$ is NOT a compact linear map}{$S$ is NOT a compact linear map and $T$ is a compact linear map}{both $S$ and $T$ are compact linear maps}{NEITHER $S$ NOR $T$ is a compact linear map}
\ans{(D)}

\textbf{Q.38} Let $c_{00} = \{(x_1, x_2, x_3, \ldots) : x_i \in \mathbb{R}, \, i \in \mathbb{N}, \, x_i \ne 0 \text{ only for finitely many indices } i\}$. For $(x_1, x_2, x_3, \ldots) \in c_{00}$, let $\|(x_1, x_2, x_3, \ldots)\|_\infty = \sup\{|x_i| : i \in \mathbb{N}\}$. Define $F, G : (c_{00}, \|\cdot\|_\infty) \to (c_{00}, \|\cdot\|_\infty)$ by
\begin{align*}
  F((x_1, x_2, \ldots, x_n, \ldots)) &= \left((1 + 1)x_1, \left(2 + \frac{1}{2}\right)x_2, \ldots, \left(n + \frac{1}{n}\right)x_n, \ldots\right), \\
  G((x_1, x_2, \ldots, x_n, \ldots)) &= \left(\frac{x_1}{1 + 1}, \frac{x_2}{2 + \frac{1}{2}}, \ldots, \frac{x_n}{n + \frac{1}{n}}, \ldots\right),
\end{align*}
for all $(x_1, x_2, \ldots, x_n, \ldots) \in c_{00}$. Then
\optslistA{$F$ is continuous but $G$ is NOT continuous}{$F$ is NOT continuous but $G$ is continuous}{both $F$ and $G$ are continuous}{NEITHER $F$ NOR $G$ is continuous}
\ans{(B)}

\textbf{Q.39} Consider the Cauchy problem
\[
  x \frac{\partial u}{\partial x} + y \frac{\partial u}{\partial y} = u; \quad u = f(t) \text{ on the initial curve } \Gamma = (t, t), \; t > 0.
\]
Consider the following statements:\\
$P$: If $f(t) = 2t + 1$, then there exists a unique solution to the Cauchy problem in a neighbourhood of $\Gamma$.\\
$Q$: If $f(t) = 2t - 1$, then there exist infinitely many solutions to the Cauchy problem in a neighbourhood of $\Gamma$.\\
Then
\optslistA{both $P$ and $Q$ are TRUE}{$P$ is FALSE and $Q$ is TRUE}{$P$ is TRUE and $Q$ is FALSE}{both $P$ and $Q$ are FALSE}
\ans{(D)}

\textbf{Q.40} Consider the linear system $Mx = b$, where $M = \begin{bmatrix} 2 & -1 \\ -4 & 3 \end{bmatrix}$ and $b = \begin{bmatrix} -2 \\ 5 \end{bmatrix}$. Suppose $M = LU$, where $L$ and $U$ are lower triangular and upper triangular square matrices, respectively. Consider the following statements:\\
$P$: If each element of the main diagonal of $L$ is $1$, then $\operatorname{trace}(U) = 3$.\\
$Q$: For any choice of the initial vector $x^{(0)}$, the Jacobi iterates $x^{(k)}$, $k = 1, 2, 3, \ldots$, converge to the unique solution of the linear system $Mx = b$.\\
Then
\optslistA{both $P$ and $Q$ are TRUE}{$P$ is FALSE and $Q$ is TRUE}{$P$ is TRUE and $Q$ is FALSE}{both $P$ and $Q$ are FALSE}
\ans{(A)}

\textbf{Q.41} Let $\phi$ and $\psi$ be two linearly independent solutions of the ordinary differential equation
\[
  y'' + (2 - \cos x)y = 0, \quad x \in \mathbb{R}.
\]
Let $\alpha, \beta \in \mathbb{R}$ be such that $\alpha < \beta$, $\phi(\alpha) = \phi(\beta) = 0$ and $\phi(x) \ne 0$ for all $x \in (\alpha, \beta)$. Consider the following statements:\\
$P$: $\phi'(\alpha)\phi'(\beta) > 0$.\\
$Q$: $\phi(x)\psi(x) \ne 0$ for all $x \in (\alpha, \beta)$.\\
Then
\optslistA{$P$ is TRUE and $Q$ is FALSE}{$P$ is FALSE and $Q$ is TRUE}{both $P$ and $Q$ are FALSE}{both $P$ and $Q$ are TRUE}
\ans{(C)}

\subsection*{Q.42 -- Q.54 Multiple Select Question (MSQ), carry TWO marks each (no negative marking).}

\textbf{Q.42 (MSQ)} Let $\mathbb{D} = \{z \in \mathbb{C} : |z| < 1\}$ and $f : \mathbb{D} \to \mathbb{C}$ be an analytic function given by the power series $f(z) = \sum_{n=0}^{\infty} a_n z^n$, where $a_0 = a_1 = 1$ and $a_n = \frac{1}{2^{2n}}$ for $n \ge 2$. Consider the following statements:\\
$P$: If $z_0 \in \mathbb{D}$, then $f$ is one-one in some neighbourhood of $z_0$.\\
$Q$: If $E = \{z \in \mathbb{C} : |z| \le \frac{1}{2}\}$, then $f(E)$ is a closed subset of $\mathbb{C}$.\\
Which of the following statements is/are correct?
\optsfourA{$P$ is TRUE}{$Q$ is TRUE}{$Q$ is FALSE}{$P$ is FALSE}
\ans{(A), (B) (MSQ)}

\textbf{Q.43 (MSQ)} Let $\Omega$ be an open connected subset of $\mathbb{C}$ containing $U = \{z \in \mathbb{C} : |z| \le \frac{1}{2}\}$. Let
\[
  \mathcal{J} = \left\{f : \Omega \to \mathbb{C} : f \text{ is analytic and } \sup_{z, w \in U} |f(z) - f(w)| = 1\right\}.
\]
Consider the following statements:\\
$P$: There exists $f \in \mathcal{J}$ such that $|f'(0)| \ge 2$.\\
$Q$: $|f^{(3)}(0)| \le 48$ for all $f \in \mathcal{J}$, where $f^{(3)}$ denotes the third derivative of $f$.\\
Then
\optsfourA{$P$ is TRUE}{$Q$ is FALSE}{$P$ is FALSE}{$Q$ is TRUE}
\ans{(C), (D) (MSQ)}

\textbf{Q.44 (MSQ)} Let $(\mathbb{R}, \tau)$ be a topological space, where the topology $\tau$ is defined as
\[
  \tau = \{U \subseteq \mathbb{R} : U = \emptyset \text{ or } 1 \in U\}.
\]
Which of the following statements is/are correct?
\optslistA{$(\mathbb{R}, \tau)$ is first countable}{$(\mathbb{R}, \tau)$ is Hausdorff}{$(\mathbb{R}, \tau)$ is separable}{The closure of $(1, 5)$ is $[1, 5]$}
\ans{(A), (C) (MSQ)}

\textbf{Q.45 (MSQ)} Let $\mathcal{R} = \{p(x) \in \mathbb{Q}[x] : p(0) \in \mathbb{Z}\}$, where $\mathbb{Q}$ denotes the set of rational numbers and $\mathbb{Z}$ denotes the set of integers. For $a \in \mathcal{R}$, let $\langle a \rangle$ denote the ideal generated by $a$ in $\mathcal{R}$. Which of the following statements is/are correct?
\optslistA{If $p(x)$ is an irreducible element in $\mathcal{R}$, then $\langle p(x) \rangle$ is a prime ideal in $\mathcal{R}$}{$\mathcal{R}$ is a unique factorization domain}{$\langle x \rangle$ is a prime ideal in $\mathcal{R}$}{$\mathcal{R}$ is NOT a principal ideal domain}
\ans{(A) (MSQ)}

\textbf{Q.46 (MSQ)} Consider the rings
\[
  \mathcal{S}_1 = \mathbb{Z}[x]/\langle 2, x^3 \rangle \quad \text{and} \quad \mathcal{S}_2 = \mathbb{Z}_2[x]/\langle x^2 \rangle,
\]
where $\langle 2, x^3 \rangle$ denotes the ideal generated by $\{2, x^3\}$ in $\mathbb{Z}[x]$ and $\langle x^2 \rangle$ denotes the ideal generated by $x^2$ in $\mathbb{Z}_2[x]$. Which of the following statements is/are correct?
\optslistA{Every prime ideal of $\mathcal{S}_1$ is a maximal ideal}{$\mathcal{S}_2$ has exactly one maximal ideal}{Every element of $\mathcal{S}_1$ is either nilpotent or a unit}{There exists an element in $\mathcal{S}_2$ which is NEITHER nilpotent NOR a unit}
\ans{(A), (B), (C) (MSQ)}

\textbf{Q.47 (MSQ)} Consider the sequence of Lebesgue measurable functions $f_n : \mathbb{R} \to \mathbb{R}$ given by
\[
  f_n(x) = \begin{cases} n^2(x - n), & \text{if } x \in \left[n, n + \frac{1}{n^2}\right], \\ 0, & \text{otherwise.} \end{cases}
\]
For a measurable subset $E$ of $\mathbb{R}$, denote $m(E)$ to be the Lebesgue measure of $E$. Which of the following statements is/are correct?
\optslistA{$\sup_{x \in \mathbb{R}} |f_n(x)| \to 0$ as $n \to \infty$}{$\int_{\mathbb{R}} |f_n(x)|\,dx \to 0$ as $n \to \infty$}{$m\!\left(\left\{x \in \mathbb{R} : |f_n(x)| > \frac{1}{2}\right\}\right) \to 0$ as $n \to \infty$}{$m(\{x \in \mathbb{R} : |f_n(x)| > 0\}) \to 0$ as $n \to \infty$}
\ans{(B), (C), (D) (MSQ)}

\textbf{Q.48 (MSQ)} Define the characteristic function $\chi_E$ of a subset $E$ in $\mathbb{R}$ by
\[
  \chi_E(x) = \begin{cases} 1, & \text{if } x \in E, \\ 0, & \text{if } x \notin E. \end{cases}
\]
For $1 \le p < 2$, let $L^p[0, 1] = \left\{f : [0, 1] \to \mathbb{R} : f \text{ is Lebesgue measurable and } \int_0^1 |f(x)|^p\,dx < \infty\right\}$. Let $f : [0, 1] \to \mathbb{R}$ be defined by
\[
  f(x) = \sum_{n=1}^{\infty} \frac{2^n}{n^3} \chi_{\left[\frac{1}{2^{n+1}}, \frac{1}{2^n}\right]}(x).
\]
Consider the following two statements:\\
$P$: $f \in L^p[0, 1]$ for every $p \in (1, 2)$.\\
$Q$: $f \in L^1[0, 1]$.\\
Then
\optsfourA{$P$ is TRUE}{$Q$ is TRUE}{$Q$ is FALSE}{$P$ is FALSE}
\ans{(B), (D) (MSQ)}

\textbf{Q.49 (MSQ)} Let $x(t), y(t)$, $t \in \mathbb{R}$, be two functions satisfying the following system of differential equations:
\[
  x'(t) = y(t), \quad y'(t) = x(t),
\]
and $x(0) = \alpha$, $y(0) = \beta$, where $\alpha, \beta$ are real numbers. Which of the following statements is/are correct?
\optslistA{If $\alpha = 1, \beta = -1$, then $|x(t)| + |y(t)| \to 0$ as $t \to \infty$}{If $\alpha = 1, \beta = 1$, then $|x(t)| + |y(t)| \to 0$ as $t \to \infty$}{If $\alpha = 1.01, \beta = -1$, then $|x(t)| + |y(t)| \to 0$ as $t \to \infty$}{If $\alpha = 1, \beta = 1.01$, then $|x(t)| + |y(t)| \to 0$ as $t \to \infty$}
\ans{(A) (MSQ)}

\textbf{Q.50 (MSQ)} For $h > 0$, and $\alpha, \beta, \gamma \in \mathbb{R}$, let
\[
  D_h f(a) = \frac{\alpha f(a - h) + \beta f(a) + \gamma f(a + 2h)}{6h}
\]
be a three-point formula to approximate $f'(a)$ for any differentiable function $f : \mathbb{R} \to \mathbb{R}$ and $a \in \mathbb{R}$. If $D_h f(a) = f'(a)$ for every polynomial $f$ of degree less than or equal to $2$ and for all $a \in \mathbb{R}$, then
\optsfourA{$\alpha + 2\gamma = -2$}{$\alpha + 2\beta - 2\gamma = 0$}{$\alpha + 2\gamma = 2$}{$\alpha + 2\beta - 2\gamma = 1$}
\ans{(A), (B) (MSQ)}

\textbf{Q.51 (MSQ)} Let $f$ be a twice continuously differentiable function on $[a, b]$ such that $f'(x) < 0$ and $f''(x) < 0$ for all $x \in (a, b)$. Let $f(\zeta) = 0$ for some $\zeta \in (a, b)$. The Newton-Raphson method to compute $\zeta$ is given by
\[
  x_{k+1} = x_k - \frac{f(x_k)}{f'(x_k)}, \quad k = 0, 1, 2, \ldots
\]
for an initial guess $x_0$. If $x_k \in (\zeta, b)$ for some $k \ge 0$, then which of the following statements is/are correct?
\optsfourA{$x_{k+1} > \zeta$}{$x_{k+1} < \zeta$}{$x_{k+1} < x_k$}{For every $\eta \in (\zeta, x_k)$, $\dfrac{f'(\eta)}{f'(x_k)} > 1$}
\ans{(A), (C) (MSQ)}

\textbf{Q.52 (MSQ)} Let $f : \mathbb{R}^2 \to \mathbb{R}$ be defined by
\[
  f(x, y) = \begin{cases} \dfrac{2x^2 y}{x^2 + y^2}, & (x, y) \ne (0, 0), \\ 0, & (x, y) = (0, 0). \end{cases}
\]
Then
\optslistA{the directional derivative of $f$ at $(0, 0)$ in the direction of $\left(\frac{1}{\sqrt{2}}, \frac{1}{\sqrt{2}}\right)$ is $\frac{1}{\sqrt{2}}$}{the directional derivative of $f$ at $(0, 0)$ in the direction of $(0, 1)$ is $1$}{the directional derivative of $f$ at $(0, 0)$ in the direction of $(1, 0)$ is $0$}{$f$ is NOT differentiable at $(0, 0)$}
\ans{(A), (C), (D) (MSQ)}

\textbf{Q.53 (MSQ)} Let $C[0, 1] = \{f : [0, 1] \to \mathbb{R} : f \text{ is continuous}\}$ and $d_\infty(f, g) = \sup\{|f(x) - g(x)| : x \in [0, 1]\}$ for $f, g \in C[0, 1]$. For each $n \in \mathbb{N}$, define $f_n : [0, 1] \to \mathbb{R}$ by $f_n(x) = x^n$ for all $x \in [0, 1]$. Let $P = \{f_n : n \in \mathbb{N}\}$. Which of the following statements is/are correct?
\optslistA{$P$ is totally bounded in $(C[0, 1], d_\infty)$}{$P$ is bounded in $(C[0, 1], d_\infty)$}{$P$ is closed in $(C[0, 1], d_\infty)$}{$P$ is open in $(C[0, 1], d_\infty)$}
\ans{(B), (C) (MSQ)}

\textbf{Q.54 (MSQ)} Let $G$ be an abelian group and $\Phi : G \to (\mathbb{Z}, +)$ be a surjective group homomorphism. Let $1 = \Phi(a)$ for some $a \in G$. Consider the following statements:\\
$P$: For every $g \in G$, there exists an $n \in \mathbb{Z}$ such that $g a^n \in \ker(\Phi)$.\\
$Q$: Let $e$ be the identity of $G$ and $\langle a \rangle$ be the subgroup generated by $a$. Then $G = \ker(\Phi)\langle a \rangle$ and $\ker(\Phi) \cap \langle a \rangle = \{e\}$.\\
Which of the following statements is/are correct?
\optsfourA{$P$ is TRUE}{$P$ is FALSE}{$Q$ is TRUE}{$Q$ is FALSE}
\ans{(A), (C) (MSQ)}

\subsection*{Q.55 -- Q.65 Numerical Answer Type (NAT), carry TWO marks each (no negative marking).}

\textbf{Q.55} Let $C$ be the curve of intersection of the cylinder $x^2 + y^2 = 4$ and the plane $z - 2 = 0$. Suppose $C$ is oriented in the counterclockwise direction around the $z$-axis, when viewed from above. If
\[
  \left|\int_C (\sin x + e^x)\,dx + 4x\,dy + e^z \cos^2 z\,dz\right| = \alpha\pi,
\]
then the value of $\alpha$ equals \underline{\hspace{2cm}}.
\ans{$16$}

\textbf{Q.56} Let $\ell^2 = \{(x_1, x_2, x_3, \ldots) : x_n \in \mathbb{R} \text{ for all } n \in \mathbb{N} \text{ and } \sum_{n=1}^{\infty} x_n^2 < \infty\}$. For a sequence $(x_1, x_2, x_3, \ldots) \in \ell^2$, define $\|(x_1, x_2, x_3, \ldots)\|_2 = \left(\sum_{n=1}^{\infty} x_n^2\right)^{1/2}$. Consider the subspace $M = \{(x_1, x_2, x_3, \ldots) \in \ell^2 : \sum_{n=1}^{\infty} \frac{x_n}{4^n} = 0\}$. Let $M^\perp$ denote the orthogonal complement of $M$ in the Hilbert space $(\ell^2, \|\cdot\|_2)$. Consider $\left(1, \frac{1}{2}, \frac{1}{3}, \frac{1}{4}, \ldots\right) \in \ell^2$. If the orthogonal projection of $\left(1, \frac{1}{2}, \frac{1}{3}, \frac{1}{4}, \ldots\right)$ onto $M^\perp$ is given by
\[
  \alpha \left(\sum_{n=1}^{\infty} \frac{1}{n 4^n}\right)\left(\frac{1}{4}, \frac{1}{4^2}, \frac{1}{4^3}, \ldots\right)
\]
for some $\alpha \in \mathbb{R}$, then $\alpha$ equals \underline{\hspace{2cm}}.
\ans{$15$}

\textbf{Q.57} Consider the transportation problem between five sources and four destinations as given in the cost table below. The supply and demand at each of the source and destination are also provided:
\begin{center}
\begin{tabular}{cccccc}
\toprule
& \multicolumn{4}{c}{\textbf{DESTINATIONS}} & \\
\cmidrule(lr){2-5}
\textbf{SOURCES} & P & Q & R & S & \textbf{Supply} \\
\midrule
1 & $13$ & $8$ & $12$ & $9$ & $20$ \\
2 & $10$ & $7$ & $5$ & $20$ & $10$ \\
3 & $3$ & $19$ & $5$ & $12$ & $50$ \\
4 & $4$ & $9$ & $7$ & $15$ & $30$ \\
5 & $14$ & $0$ & $1$ & $7$ & $40$ \\
\midrule
\textbf{Demand} & $60$ & $10$ & $20$ & $60$ & \\
\bottomrule
\end{tabular}
\end{center}
Let $C_N$ and $C_L$ be the total cost of the initial basic feasible solution obtained from the North-West corner method and the Least-Cost method, respectively. Then $C_N - C_L$ equals \underline{\hspace{2cm}}.
\ans{$380$}

\textbf{Q.58} Let $\sigma \in S_8$, where $S_8$ is the permutation group on $8$ elements. Suppose $\sigma$ is the product of $\sigma_1$ and $\sigma_2$, where $\sigma_1$ is a $4$-cycle and $\sigma_2$ is a $3$-cycle in $S_8$. If $\sigma_1$ and $\sigma_2$ are disjoint cycles, then the number of elements in $S_8$ which are conjugate to $\sigma$ is \underline{\hspace{2cm}}.
\ans{$3360$}

\textbf{Q.59} Let $A$ be a $3 \times 3$ real matrix with $\det(A + iI) = 0$, where $i = \sqrt{-1}$ and $I$ is the $3 \times 3$ identity matrix. If $\det(A) = 3$, then the trace of $A^2$ is \underline{\hspace{2cm}}.
\ans{$7$}

\textbf{Q.60} Let $A = [a_{ij}]$ be a $3 \times 3$ real matrix such that
\[
  A\begin{bmatrix} 1 \\ 2 \\ 1 \end{bmatrix} = 2\begin{bmatrix} 1 \\ 2 \\ 1 \end{bmatrix}, \quad
  A\begin{bmatrix} 0 \\ 1 \\ 1 \end{bmatrix} = 2\begin{bmatrix} 0 \\ 1 \\ 1 \end{bmatrix} \quad \text{and} \quad
  A\begin{bmatrix} -1 \\ 1 \\ 0 \end{bmatrix} = 4\begin{bmatrix} -1 \\ 1 \\ 0 \end{bmatrix}.
\]
If $m$ is the degree of the minimal polynomial of $A$, then $a_{11} + a_{21} + a_{31} + m$ equals \underline{\hspace{2cm}}.
\ans{$4$}

\textbf{Q.61} Let $\Omega$ be the disk $x^2 + y^2 < 4$ in $\mathbb{R}^2$ with boundary $\partial\Omega$. If $u(x, y)$ is the solution of the Dirichlet problem
\begin{align*}
  \frac{\partial^2 u}{\partial x^2} + \frac{\partial^2 u}{\partial y^2} &= 0, \quad (x, y) \in \Omega, \\
  u(x, y) &= 1 + 2x^2, \quad (x, y) \in \partial\Omega,
\end{align*}
then the value of $u(0, 1)$ is \underline{\hspace{2cm}}.
\ans{$4$}

\textbf{Q.62} For every $k \in \mathbb{N} \cup \{0\}$, let $y_k(x)$ be a polynomial of degree $k$ with $y_k(1) = 5$. Further, let $y_k(x)$ satisfy the Legendre equation
\[
  (1 - x^2)y'' - 2xy' + k(k + 1)y = 0.
\]
If
\[
  \frac{1}{2} \int_{-1}^{1} \sum_{k=1}^{n} \bigl(y_k(x) - y_{k-1}(x)\bigr)^2\,dx - \int_{-1}^{1} \sum_{k=1}^{n} \bigl(y_k(x)\bigr)^2\,dx = 24,
\]
for some positive integer $n$, then the value of $n$ is \underline{\hspace{2cm}}.
\ans{$12$}

\textbf{Q.63} Consider the ordinary differential equation (ODE)
\[
  4(\ln x)y'' + 3y' + y = 0, \quad x > 1.
\]
If $r_1$ and $r_2$ are the roots of the indicial equation of the above ODE at the regular singular point $x = 1$, then $|r_1 - r_2|$ is equal to \underline{\hspace{2cm}} (rounded off to 2 decimal places).
\ans{$0.24 \text{ to } 0.26$}

\textbf{Q.64} Let $u(x, t)$ be the solution of the non-homogeneous wave equation
\begin{align*}
  \frac{\partial^2 u}{\partial x^2} - \frac{\partial^2 u}{\partial t^2} &= \sin x \sin(2t), \quad 0 < x < \pi, \quad t > 0, \\
  u(x, 0) &= 0, \quad \text{and} \quad \frac{\partial u}{\partial t}(x, 0) = 0, \quad \text{for } 0 \le x \le \pi, \\
  u(0, t) &= 0, \quad u(\pi, t) = 0, \quad \text{for } t \ge 0.
\end{align*}
Then the value of $u\!\left(\frac{\pi}{2}, \frac{3\pi}{2}\right)$ is \underline{\hspace{2cm}} (rounded off to 2 decimal places).
\ans{$0.64 \text{ to } 0.70$}

\textbf{Q.65} Consider the Linear Programming Problem P:
\begin{align*}
  \text{Maximize} \quad & 3x_1 + 2x_2 + 5x_3 \\
  \text{subject to} \quad & x_1 + 2x_2 + x_3 \le 44, \\
  & x_1 + 2x_3 \le 48, \\
  & x_1 + 4x_2 \le 52, \\
  & x_1 \ge 0, \quad x_2 \ge 0, \quad x_3 \ge 0.
\end{align*}
The optimal value of the problem P is equal to \underline{\hspace{2cm}}.
\ans{$140$}

% ============================================================================
% GATE 2024 -- GENERAL APTITUDE (GA) & MATHEMATICS (MA)
% ============================================================================

\newpage
\begin{center}
  {\LARGE\bfseries GATE 2024}\\[4pt]
  {\large\bfseries General Aptitude (GA) \& Mathematics (MA)}\\[2pt]
  {\small\itshape Official Final Answer Key attached to every question}
\end{center}
\medskip

\section*{General Aptitude (GA)}

\subsection*{Q.1 -- Q.5 Multiple Choice Question (MCQ), carry ONE mark each (for each wrong answer: $-1/3$).}

\textbf{Q.1} If `$\to$' denotes increasing order of intensity, then the meaning of the words $[\text{drizzle} \to \text{rain} \to \text{downpour}]$ is analogous to $[\underline{\hspace{1.5cm}} \to \text{quarrel} \to \text{feud}]$. Which one of the given options is appropriate to fill the blank?
\optsfourA{bicker}{bog}{dither}{dodge}
\ans{(A)}

\textbf{Q.2} Statements:
\begin{enumerate}
  \setlength{\itemsep}{1pt}\setlength{\parskip}{0pt}
  \item All heroes are winners.
  \item All winners are lucky people.
\end{enumerate}
Inferences:
\begin{itemize}
  \setlength{\itemsep}{1pt}\setlength{\parskip}{0pt}
  \item[I.] All lucky people are heroes.
  \item[II.] Some lucky people are heroes.
  \item[III.] Some winners are heroes.
\end{itemize}
Which of the above inferences can be logically deduced from statements 1 and 2?
\optsfourA{Only I and II}{Only II and III}{Only I and III}{Only III}
\ans{(B)}

\textbf{Q.3} A student was supposed to multiply a positive real number $p$ with another positive real number $q$. Instead, the student divided $p$ by $q$. If the percentage error in the student's answer is $80\%$, the value of $q$ is
\optsfourA{$5$}{$\sqrt{2}$}{$2$}{$\sqrt{5}$}
\ans{(D)}

\textbf{Q.4} If the sum of the first $20$ consecutive positive odd numbers is divided by $20^2$, the result is
\optsfourA{$1$}{$20$}{$2$}{$1/2$}
\ans{(A)}

\textbf{Q.5} The ratio of the number of girls to boys in class VIII is the same as the ratio of the number of boys to girls in class IX. The total number of students (boys and girls) in classes VIII and IX is $450$ and $360$, respectively. If the number of girls in classes VIII and IX is the same, then the number of girls in each class is
\optsfourA{$150$}{$200$}{$250$}{$175$}
\ans{(B)}

\subsection*{Q.6 -- Q.10 Multiple Choice Question (MCQ), carry TWO marks each (for each wrong answer: $-2/3$).}

\textbf{Q.6} In the given text, the blanks are numbered (i)--(iv). Select the best match for all the blanks.\\
Yoko Roi stands \underline{\hspace{0.7cm}(i)\hspace{0.7cm}} as an author for standing \underline{\hspace{0.7cm}(ii)\hspace{0.7cm}} as an honorary fellow, after she stood \underline{\hspace{0.7cm}(iii)\hspace{0.7cm}} her writings that stand \underline{\hspace{0.7cm}(iv)\hspace{0.7cm}} the freedom of speech.
\optslistA{(i) out \quad (ii) down \quad (iii) in \quad (iv) for}{(i) down \quad (ii) out \quad (iii) by \quad (iv) in}{(i) down \quad (ii) out \quad (iii) for \quad (iv) in}{(i) out \quad (ii) down \quad (iii) by \quad (iv) for}
\ans{(D)}

\textbf{Q.7} Seven identical cylindrical chalk-sticks are fitted tightly in a cylindrical container. The figure below shows the arrangement of the chalk-sticks inside the cylinder.\\
\emph{[Diagram note: Cross-sectional view showing one central circle of radius $r$ tightly surrounded by $6$ mutually tangent identical circles of radius $r$, all inscribed tightly inside a larger circular container of radius $3r$.]}\\
The length of the container is equal to the length of the chalk-sticks. The ratio of the occupied space to the empty space of the container is
\optsfourA{$5/2$}{$7/2$}{$9/2$}{$3$}
\ans{(B)}

\textbf{Q.8} The plot below shows the relationship between the mortality risk of cardiovascular disease and the number of steps a person walks per day. Based on the data, which one of the following options is true?\\[4pt]
\emph{[Chart data: Smoothly decreasing convex curve of Mortality Risk vs.\ Steps/Day: $1.0$ at $0$ steps/day, $\approx 0.45$ at $5000$, $\approx 0.18$ at $10000$, $\approx 0.09$ at $15000$, and $\approx 0.04$ at $20000$ steps/day.]}
\optslistA{The risk reduction on increasing the steps/day from $0$ to $10000$ is less than the risk reduction on increasing the steps/day from $10000$ to $20000$.}{The risk reduction on increasing the steps/day from $0$ to $5000$ is less than the risk reduction on increasing the steps/day from $15000$ to $20000$.}{For any $5000$ increment in steps/day the largest risk reduction occurs on going from $0$ to $5000$.}{For any $5000$ increment in steps/day the largest risk reduction occurs on going from $15000$ to $20000$.}
\ans{(C)}

\textbf{Q.9} Five cubes of identical size and another smaller cube are assembled as shown in Figure A. If viewed from direction X, the planar image of the assembly appears as Figure B.\\
\emph{[Diagram note: Figure A shows a $2 \times 2$ base layer of four identical large cubes, with a fifth identical large cube stacked on the back-left base cube and a smaller half-size cube placed on the outer-front corner of the back-right base cube. Viewed from above (direction X), Figure B is a $2 \times 2$ grid of squares with the small square nested in the bottom-right corner of the bottom-right cell. Direction Y looks horizontally from the front-right side toward the back-left.]}\\
If viewed from direction Y, the planar image of the assembly (Figure A) will appear as
\optslistA{Two base squares in the bottom row, one large square on the top-right with the small square nested in its bottom-right corner}{Two base squares in the bottom row, one large square on the top-right divided vertically in half}{Two base squares in the bottom row, one large square on the top-right with the small square nested in its bottom-left corner}{Two base squares in the bottom row, one large square on the top-right with the small square sitting on top of the bottom-left base square}
\ans{(A)}

\textbf{Q.10} Visualize a cube that is held with one of the four body diagonals aligned to the vertical axis. Rotate the cube about this axis such that its view remains unchanged. The magnitude of the minimum angle of rotation is
\optsfourA{$120^\circ$}{$60^\circ$}{$90^\circ$}{$180^\circ$}
\ans{(A)}

\section*{Mathematics (MA)}

\textbf{Useful Data:} $\mathbb{Z}$: set of integers; $\mathbb{R}$: set of real numbers; $\mathbb{C}$: set of complex numbers; $\operatorname{Im}(z)$: imaginary part of $z$; $\operatorname{Re}(z)$: real part of $z$; $\sup(A)$, $\inf(A)$: supremum and infimum of $A \subseteq \mathbb{R}$; $M_n(\mathbb{F})$: $n \times n$ matrices over field $\mathbb{F}$; $GL_n(\mathbb{F})$: $n \times n$ invertible matrices over $\mathbb{F}$; $\mathbb{F}^n$: $n$-dimensional vector space over $\mathbb{F}$; $I_n$: $n \times n$ identity matrix; $S_n$: symmetric group on $n$ elements; $A \setminus B = \{a \in A : a \notin B\}$; $\mathbb{Z}/n\mathbb{Z}$: cyclic group of order $n$; $\ln$: natural logarithm; $f', f'', f'''$: first-, second-, and third-order derivatives of a one-variable function $f$.

\subsection*{Q.11 -- Q.25 Multiple Choice Question (MCQ), carry ONE mark each (for each wrong answer: $-1/3$).}

\textbf{Q.11} Consider the following condition on a function $f : \mathbb{C} \to \mathbb{C}$:
\[
  |f(z)| = 1 \quad \text{for all } z \in \mathbb{C} \text{ such that } \operatorname{Im}(z) = 0. \tag{P}
\]
Which one of the following is correct?
\optslistA{There is a non-constant analytic polynomial $f$ satisfying (P)}{Every entire function $f$ satisfying (P) is a constant function}{Every entire function $f$ satisfying (P) has no zeroes in $\mathbb{C}$}{There is an entire function $f$ satisfying (P) with infinitely many zeroes in $\mathbb{C}$}
\ans{(C)}

\textbf{Q.12} Let $C$ be the ellipse $\{z \in \mathbb{C} : |z - 2| + |z + 2| = 8\}$ traversed counter-clockwise. The value of the contour integral
\[
  \oint_C \frac{z^2\,dz}{z^2 - 2z + 2}
\]
is equal to
\optsfourA{$0$}{$2\pi i$}{$4\pi i$}{$-\pi i$}
\ans{(C)}

\textbf{Q.13} Let $X$ be a topological space and $A \subseteq X$. Given a subset $S$ of $X$, let $\operatorname{int}(S)$, $\partial S$, and $\overline{S}$ denote the interior, boundary, and closure, respectively, of the set $S$. Which one of the following is NOT necessarily true?
\optsfourA{$\operatorname{int}(X \setminus A) \subseteq X \setminus \overline{A}$}{$A \subseteq \overline{A}$}{$\partial A \subseteq \partial(\operatorname{int}(A))$}{$\partial(\overline{A}) \subseteq \partial A$}
\ans{(C)}

\textbf{Q.14} Consider the following limit
\[
  \lim_{\varepsilon \to 0} \frac{1}{\varepsilon} \int_0^\infty e^{-x/\varepsilon} \bigl(\cos(3x) + x^2 + \sqrt{x + 4}\bigr)\,dx.
\]
Which one of the following is correct?
\optslistA{The limit does not exist}{The limit exists and is equal to $0$}{The limit exists and is equal to $3$}{The limit exists and is equal to $\pi$}
\ans{(C)}

\textbf{Q.15} Let $\mathbb{R}[X^2, X^3]$ be the subring of $\mathbb{R}[X]$ generated by $X^2$ and $X^3$. Consider the following statements:\\
I. The ring $\mathbb{R}[X^2, X^3]$ is a unique factorization domain.\\
II. The ring $\mathbb{R}[X^2, X^3]$ is a principal ideal domain.\\
Which one of the following is correct?
\optslistA{Both I and II are TRUE}{I is TRUE and II is FALSE}{I is FALSE and II is TRUE}{Both I and II are FALSE}
\ans{(D)}

\textbf{Q.16} Given a prime number $p$, let $n_p(G)$ denote the number of $p$-Sylow subgroups of a finite group $G$. Which one of the following is TRUE for every group $G$ of order $2024$?
\optslistA{$n_{11}(G) = 1$ and $n_{23}(G) = 11$}{$n_{11}(G) \in \{1, 23\}$ and $n_{23}(G) = 1$}{$n_{11}(G) = 23$ and $n_{23}(G) \in \{1, 88\}$}{$n_{11}(G) = 23$ and $n_{23}(G) = 11$}
\ans{(B)}

\textbf{Q.17} Consider the following statements:\\
I. Every compact Hausdorff space is normal.\\
II. Every metric space is normal.\\
Which one of the following is correct?
\optslistA{Both I and II are TRUE}{I is TRUE and II is FALSE}{I is FALSE and II is TRUE}{Both I and II are FALSE}
\ans{(A)}

\textbf{Q.18} Consider the topology on $\mathbb{Z}$ with basis $\{S(a, b) : a, b \in \mathbb{Z} \text{ and } a \ne 0\}$, where $S(a, b) = \{an + b : n \in \mathbb{Z}\}$. Consider the following statements:\\
I. $S(a, b)$ is both open and closed for each $a, b \in \mathbb{Z}$ with $a \ne 0$.\\
II. The only connected set containing $x \in \mathbb{Z}$ is $\{x\}$.\\
Which one of the following is correct?
\optslistA{Both I and II are TRUE}{I is TRUE and II is FALSE}{I is FALSE and II is TRUE}{Both I and II are FALSE}
\ans{(A)}

\textbf{Q.19} Let $A = \begin{pmatrix} 0 & 2 \\ 2 & 0 \end{pmatrix}$ and $T : M_2(\mathbb{C}) \to M_2(\mathbb{C})$ be the linear transformation given by $T(B) = AB$. The characteristic polynomial of $T$ is
\optsfourA{$X^4 - 8X^2 + 16$}{$X^2 - 4$}{$X^2 - 2$}{$X^4 - 16$}
\ans{(A)}

\textbf{Q.20} Let $A \in M_n(\mathbb{C})$ be a normal matrix. Consider the following statements:\\
I. If all the eigenvalues of $A$ are real, then $A$ is Hermitian.\\
II. If all the eigenvalues of $A$ have absolute value $1$, then $A$ is unitary.\\
Which one of the following is correct?
\optslistA{Both I and II are TRUE}{I is TRUE and II is FALSE}{I is FALSE and II is TRUE}{Both I and II are FALSE}
\ans{(A)}

\textbf{Q.21} Let
\[
  A = \begin{pmatrix} 2 & -1 & 1 \\ 1 & 2 & -1 \\ -1 & 1 & 2 \end{pmatrix}
\]
and $b$ be a $3 \times 1$ real column vector. Consider the statements:\\
I. The Jacobi iteration method for the system $(A + \varepsilon I_3)x = b$ converges for any initial approximation and $\varepsilon > 0$.\\
II. The Gauss--Seidel iteration method for the system $(A + \varepsilon I_3)x = b$ converges for any initial approximation and $\varepsilon > 0$.\\
Which one of the following is correct?
\optslistA{Both I and II are TRUE}{I is TRUE and II is FALSE}{I is FALSE and II is TRUE}{Both I and II are FALSE}
\ans{(A)}

\textbf{Q.22} For the initial value problem
\[
  y' = f(x, y), \quad y(x_0) = y_0,
\]
generate approximations $y_n$ to $y(x_n)$, $x_n = x_0 + nh$, for a fixed $h > 0$ and $n = 1, 2, 3, \ldots$, using the recursion formula
\[
  y_n = y_{n-1} + a k_1 + b k_2, \quad \text{where } k_1 = h f(x_{n-1}, y_{n-1}) \text{ and } k_2 = h f(x_{n-1} + \alpha h, y_{n-1} + \beta k_1).
\]
Which one of the following choices of $a, b, \alpha, \beta$ for the above recursion formula gives the Runge--Kutta method of order $2$?
\optslistA{$a = 1, \; b = 1, \; \alpha = 0.5, \; \beta = 0.5$}{$a = 0.5, \; b = 0.5, \; \alpha = 2, \; \beta = 2$}{$a = 0.25, \; b = 0.75, \; \alpha = 2/3, \; \beta = 2/3$}{$a = 0.5, \; b = 0.5, \; \alpha = 1, \; \beta = 2$}
\ans{(C)}

\textbf{Q.23} Let $u = u(x, t)$ be the solution of
\begin{align*}
  \frac{\partial u}{\partial t} - 4 \frac{\partial^2 u}{\partial x^2} &= 0, \quad 0 < x < 1, \; t > 0, \\
  u(0, t) &= u(1, t) = 0, \quad t \ge 0, \\
  u(x, 0) &= \sin(\pi x), \quad 0 \le x \le 1.
\end{align*}
Define $g(t) = \int_0^1 (u(x, t))^2\,dx$, for $t > 0$. Which one of the following is correct?
\optslistA{$g$ is decreasing on $(0, \infty)$ and $\lim_{t \to \infty} g(t) = 0$}{$g$ is decreasing on $(0, \infty)$ and $\lim_{t \to \infty} g(t) = 1/4$}{$g$ is increasing on $(0, \infty)$ and $\lim_{t \to \infty} g(t)$ does not exist}{$g$ is increasing on $(0, \infty)$ and $\lim_{t \to \infty} g(t) = 3$}
\ans{(A)}

\textbf{Q.24} If $y_1$ and $y_2$ are two different solutions of the ordinary differential equation
\[
  y' + \sin(e^x) y = \cos(e^{x+1}), \quad 0 \le x \le 1,
\]
then which one of the following is its general solution on $[0, 1]$?
\optslistA{$c_1 y_1 + c_2 y_2, \; c_1, c_2 \in \mathbb{R}$}{$y_1 + c(y_1 - y_2), \; c \in \mathbb{R}$}{$c y_1 + (y_1 - y_2), \; c \in \mathbb{R}$}{$c_1(y_1 + y_2) + c_2(y_1 - y_2), \; c_1, c_2 \in \mathbb{R}$}
\ans{(B)}

\textbf{Q.25} Consider the following Linear Programming Problem P:
\begin{align*}
  \text{minimize} \quad & 5x_1 + 2x_2 \\
  \text{subject to} \quad & 2x_1 + x_2 \le 2, \\
  & x_1 + x_2 \ge 1, \\
  & x_1, x_2 \ge 0.
\end{align*}
The optimal value of the problem P is equal to
\optsfourA{$5$}{$0$}{$4$}{$2$}
\ans{(D)}

\subsection*{Q.26 -- Q.35 Numerical Answer Type (NAT), carry ONE mark each (no negative marking).}

\textbf{Q.26} Let $p = \left(1, \frac{1}{2}, \frac{1}{3}, \frac{1}{4}\right) \in \mathbb{R}^4$ and $f : \mathbb{R}^4 \to \mathbb{R}$ be a differentiable function such that $f(p) = 6$ and $f(\lambda x) = \lambda^3 f(x)$, for every $\lambda \in (0, \infty)$ and $x \in \mathbb{R}^4$. The value of
\[
  12 \frac{\partial f}{\partial x_1}(p) + 6 \frac{\partial f}{\partial x_2}(p) + 4 \frac{\partial f}{\partial x_3}(p) + 3 \frac{\partial f}{\partial x_4}(p)
\]
is equal to \underline{\hspace{2cm}} (answer in integer).
\ans{$216$}

\textbf{Q.27} The number of non-isomorphic finite groups with exactly $3$ conjugacy classes is equal to \underline{\hspace{2cm}} (answer in integer).
\ans{$2$}

\textbf{Q.28} Let $f(x, y) = (x^2 - y^2, 2xy)$, where $x > 0, y > 0$. Let $g$ be the inverse of $f$ in a neighborhood of $f(2, 1)$. Then the determinant of the Jacobian matrix of $g$ at $f(2, 1)$ is equal to \underline{\hspace{2cm}} (round off to TWO decimal places).
\ans{$0.04 \text{ to } 0.06$}

\textbf{Q.29} Let $\mathbb{F}_3$ be the field with exactly $3$ elements. The number of elements in $GL_2(\mathbb{F}_3)$ is equal to \underline{\hspace{2cm}} (answer in integer).
\ans{$48$}

\textbf{Q.30} Given a real subspace $W$ of $\mathbb{R}^4$, let $W^\perp$ denote its orthogonal complement with respect to the standard inner product on $\mathbb{R}^4$. Let $W_1 = \operatorname{Span}\{(1, 0, 0, -1)\}$ and $W_2 = \operatorname{Span}\{(2, 1, 0, -1)\}$ be real subspaces of $\mathbb{R}^4$. The dimension of $W_1^\perp \cap W_2^\perp$ over $\mathbb{R}$ is equal to \underline{\hspace{2cm}} (answer in integer).
\ans{$2$}

\textbf{Q.31} The number of group homomorphisms from $\mathbb{Z}/4\mathbb{Z}$ to $S_4$ is equal to \underline{\hspace{2cm}} (answer in integer).
\ans{$16$}

\textbf{Q.32} Let $a \in \mathbb{R}$ and $h$ be a positive real number. For any twice-differentiable function $f : \mathbb{R} \to \mathbb{R}$, let $P_f(x)$ be the interpolating polynomial of degree at most two that interpolates $f$ at the points $a - h, a, a + h$. Define $d$ to be the largest integer such that any polynomial $g$ of degree $d$ satisfies
\[
  g''(a) = P_g''(a).
\]
The value of $d$ is equal to \underline{\hspace{2cm}} (answer in integer).
\ans{$3$}

\textbf{Q.33} Let $P_f(x)$ be the interpolating polynomial of degree at most two that interpolates the function $f(x) = x^2|x|$ at the points $x = -1, 0, 1$. Then
\[
  \sup_{x \in [-1, 1]} |f(x) - P_f(x)| = \underline{\hspace{2cm}} \quad \text{(round off to TWO decimal places)}.
\]
\ans{$0.14 \text{ to } 0.16$}

\textbf{Q.34} The maximum of the function $f(x, y, z) = xyz$ subject to the constraints
\[
  xy + yz + zx = 12, \quad x > 0, \; y > 0, \; z > 0,
\]
is equal to \underline{\hspace{2cm}} (round off to TWO decimal places).
\ans{$7.99 \text{ to } 8.01$}

\textbf{Q.35} If the outward flux of $\mathbf{F}(x, y, z) = (x^3, y^3, z^3)$ through the unit sphere $x^2 + y^2 + z^2 = 1$ is $\alpha\pi$, then $\alpha$ is equal to \underline{\hspace{2cm}} (round off to TWO decimal places).
\ans{$2.39 \text{ to } 2.41$}

\subsection*{Q.36 -- Q.47 Multiple Choice Question (MCQ), carry TWO marks each (for each wrong answer: $-2/3$).}

\textbf{Q.36} Let $\mathbb{H} = \{z \in \mathbb{C} : \operatorname{Im}(z) > 0\}$ and $\mathbb{D} = \{z \in \mathbb{C} : |z| < 1\}$. Then
\[
  \sup\left\{|f'(0)| : f \text{ is an analytic function from } \mathbb{D} \text{ to } \mathbb{H} \text{ and } f(0) = \frac{i}{2}\right\}
\]
is equal to
\optsfourA{$\dfrac{1}{4}$}{$\dfrac{1}{2}$}{$1$}{$100$}
\ans{(C)}

\textbf{Q.37} Let $S^1 = \{z \in \mathbb{C} : |z| = 1\}$. For which one of the following functions $f$ does there exist a sequence of polynomials in $z$ that uniformly converges to $f$ on $S^1$?
\optsfourA{$f(z) = \overline{z}$}{$f(z) = \operatorname{Re}(z)$}{$f(z) = e^z$}{$f(z) = |z + 1|^2$}
\ans{(C)}

\textbf{Q.38} Let $f : [0, 1] \to \mathbb{R}$ be a function. Which one of the following is a sufficient condition for $f$ to be Lebesgue measurable?
\optslistA{$|f|$ is a Lebesgue measurable function}{There exist continuous functions $g, h : [0, 1] \to \mathbb{R}$ such that $g \le f \le h$ on $[0, 1]$}{$f$ is continuous almost everywhere on $[0, 1]$}{For each $c \in \mathbb{R}$, the set $\{x \in [0, 1] : f(x) = c\}$ is Lebesgue measurable}
\ans{(C)}

\textbf{Q.39} Let $g : M_2(\mathbb{R}) \to \mathbb{R}$ be given by $g(A) = \operatorname{Trace}(A^2)$. Let $0$ be the $2 \times 2$ zero matrix. The space $M_2(\mathbb{R})$ may be identified with $\mathbb{R}^4$ in the usual manner. Which one of the following is correct?
\optslistA{$0$ is a point of local minimum of $g$}{$0$ is a point of local maximum of $g$}{$0$ is a saddle point of $g$}{$0$ is not a critical point of $g$}
\ans{(C)}

\textbf{Q.40} Consider the following statements:\\
I. There exists a proper subgroup $G$ of $(\mathbb{Q}, +)$ such that $\mathbb{Q}/G$ is a finite group.\\
II. There exists a subgroup $G$ of $(\mathbb{Q}, +)$ such that $\mathbb{Q}/G$ is isomorphic to $(\mathbb{Z}, +)$.\\
Which one of the following is correct?
\optslistA{Both I and II are TRUE}{I is TRUE and II is FALSE}{I is FALSE and II is TRUE}{Both I and II are FALSE}
\ans{(D)}

\textbf{Q.41} Let $X$ be the space $\mathbb{R}/\mathbb{Z}$ with the quotient topology induced from the usual topology on $\mathbb{R}$. Consider the following statements:\\
I. $X$ is compact.\\
II. $X \setminus \{x\}$ is connected for any $x \in X$.\\
Which one of the following is correct?
\optslistA{Both I and II are TRUE}{I is TRUE and II is FALSE}{I is FALSE and II is TRUE}{Both I and II are FALSE}
\ans{(A)}

\textbf{Q.42} Let $\langle \cdot, \cdot \rangle$ denote the standard inner product on $\mathbb{R}^7$. Let $\Sigma = \{v_1, \ldots, v_5\} \subseteq \mathbb{R}^7$ be a set of unit vectors such that $\langle v_i, v_j \rangle$ is a non-positive integer for all $1 \le i \ne j \le 5$. Define $N(\Sigma)$ to be the number of pairs $(r, s)$, $1 \le r, s \le 5$, such that $\langle v_r, v_s \rangle \ne 0$. The maximum possible value of $N(\Sigma)$ is equal to
\optsfourA{$9$}{$10$}{$14$}{$5$}
\ans{(A)}

\textbf{Q.43} Let $f(x) = |x| + |x - 1| + |x - 2|$, $x \in [-1, 2]$. Which one of the following numerical integration rules gives the exact value of the integral
\[
  \int_{-1}^{2} f(x)\,dx?
\]
\optslistA{The Simpson's rule}{The trapezoidal rule}{The composite Simpson's rule by dividing $[-1, 2]$ into $4$ equal subintervals}{The composite trapezoidal rule by dividing $[-1, 2]$ into $3$ equal subintervals}
\ans{(D)}

\textbf{Q.44} Consider the initial value problem (IVP)
\[
  y' = e^{-y^2} + 1, \quad y(0) = 0.
\]
I. IVP has a unique solution on $\mathbb{R}$.\\
II. Every solution of IVP is bounded on its maximal interval of existence.\\
Which one of the following is correct?
\optslistA{Both I and II are TRUE}{I is TRUE and II is FALSE}{I is FALSE and II is TRUE}{Both I and II are FALSE}
\ans{(B)}

\textbf{Q.45} Let $A$ be a $2 \times 2$ non-diagonalizable real matrix with a real eigenvalue $\lambda$ and $v$ be an eigenvector of $A$ corresponding to $\lambda$. Which one of the following is the general solution of the system $y' = Ay$ of first-order linear differential equations?
\optslistA{$c_1 e^{\lambda t} v + c_2 t e^{\lambda t} v$, where $c_1, c_2 \in \mathbb{R}$}{$c_1 e^{\lambda t} v + c_2 t^2 e^{\lambda t} v$, where $c_1, c_2 \in \mathbb{R}$}{$c_1 e^{\lambda t} v + c_2 e^{\lambda t}(tv + u)$, where $c_1, c_2 \in \mathbb{R}$ and $u$ is a $2 \times 1$ real column vector such that $(A - \lambda I_2)u = v$}{$c_1 e^{\lambda t} v + c_2 t e^{\lambda t}(v + u)$, where $c_1, c_2 \in \mathbb{R}$ and $u$ is a $2 \times 1$ real column vector such that $(A - \lambda I_2)u = v$}
\ans{(C)}

\textbf{Q.46} Let $D = \{(x, y) \in \mathbb{R}^2 : x > 0 \text{ and } y > 0\}$. If the following second-order linear partial differential equation
\[
  y^2 \frac{\partial^2 u}{\partial x^2} - x^2 \frac{\partial^2 u}{\partial y^2} + y \frac{\partial u}{\partial y} = 0 \quad \text{on } D
\]
is transformed to
\[
  \left(\frac{\partial^2 u}{\partial \eta^2} - \frac{\partial^2 u}{\partial \xi^2}\right) + \left(\frac{\partial u}{\partial \eta} + \frac{\partial u}{\partial \xi}\right)\frac{1}{2\eta} + \left(a \frac{\partial u}{\partial \eta} + b \frac{\partial u}{\partial \xi}\right)\frac{1}{2\xi} = 0 \quad \text{on } D
\]
for some $a, b \in \mathbb{R}$, via the co-ordinate transform $\eta = \frac{x^2}{2}$ and $\xi = \frac{y^2}{2}$, then which one of the following is correct?
\optsfourA{$a = 2, \; b = 0$}{$a = 0, \; b = -1$}{$a = 1, \; b = -1$}{$a = 1, \; b = 0$}
\ans{(B)}

\textbf{Q.47} Let $\ell^p = \{x = (x_n)_{n \ge 1} : x_n \in \mathbb{R}, \, \|x\|_p = (\sum_{n=1}^{\infty} |x_n|^p)^{1/p} < \infty\}$ for $p = 1, 2$. Let $C_{00} = \{(x_n)_{n \ge 1} : x_n = 0 \text{ for all but finitely many } n \ge 1\}$. For $x = (x_n)_{n \ge 1} \in C_{00}$, define $f(x) = \sum_{n=1}^{\infty} \frac{x_n}{\sqrt{n}}$. Consider the following statements:\\
I. There exists a continuous linear functional $F$ on $(\ell^1, \|\cdot\|_1)$ such that $F = f$ on $C_{00}$.\\
II. There exists a continuous linear functional $G$ on $(\ell^2, \|\cdot\|_2)$ such that $G = f$ on $C_{00}$.\\
Which one of the following is correct?
\optslistA{Both I and II are TRUE}{I is TRUE and II is FALSE}{I is FALSE and II is TRUE}{Both I and II are FALSE}
\ans{(B)}

\subsection*{Q.48 -- Q.56 Multiple Select Question (MSQ), carry TWO marks each (no negative marking).}

\textbf{Q.48 (MSQ)} Let $\ell^2_{\mathbb{Z}} = \left\{(x_j)_{j \in \mathbb{Z}} : x_j \in \mathbb{R} \text{ and } \sum_{j=-\infty}^{\infty} x_j^2 < \infty\right\}$ endowed with the inner product
\[
  \langle x, y \rangle = \sum_{j=-\infty}^{\infty} x_j y_j, \quad x = (x_j)_{j \in \mathbb{Z}}, \; y = (y_j)_{j \in \mathbb{Z}} \in \ell^2_{\mathbb{Z}}.
\]
Let $T : \ell^2_{\mathbb{Z}} \to \ell^2_{\mathbb{Z}}$ be given by $T((x_j)_{j \in \mathbb{Z}}) = (y_j)_{j \in \mathbb{Z}}$, where
\[
  y_j = \frac{x_j + x_{-j}}{2}, \quad j \in \mathbb{Z}.
\]
Which of the following is/are correct?
\optslistA{$T$ is a compact operator}{The operator norm of $T$ is $1$}{$T$ is a self-adjoint operator}{$\operatorname{Range}(T)$ is closed}
\ans{(B), (C), (D) (MSQ)}

\textbf{Q.49 (MSQ)} Let $X$ be the normed space $(\mathbb{R}^2, \|\cdot\|)$, where $\|(x, y)\| = |x| + |y|$, $(x, y) \in \mathbb{R}^2$. Let $S = \{(x, 0) : x \in \mathbb{R}\}$ and $f : S \to \mathbb{R}$ be given by $f((x, 0)) = 2x$ for all $x \in \mathbb{R}$. Recall that a Hahn--Banach extension of $f$ to $X$ is a continuous linear functional $F$ on $X$ such that $F|_S = f$ and $\|F\| = \|f\|$, where $\|F\|$ and $\|f\|$ are the norms of $F$ and $f$ on $X$ and $S$, respectively. Which of the following is/are true?
\optslistA{$F(x, y) = 2x + 3y$ is a Hahn--Banach extension of $f$ to $X$}{$F(x, y) = 2x + y$ is a Hahn--Banach extension of $f$ to $X$}{$f$ admits infinitely many Hahn--Banach extensions to $X$}{$f$ admits exactly two distinct Hahn--Banach extensions to $X$}
\ans{(B), (C) (MSQ)}

\textbf{Q.50 (MSQ)} Let $\{[a, b) : a, b \in \mathbb{R}, \, a < b\}$ be a basis for a topology $\tau$ on $\mathbb{R}$. Which of the following is/are correct?
\optslistA{Every $(a, b)$ with $a < b$ is an open set in $(\mathbb{R}, \tau)$}{Every $[a, b]$ with $a < b$ is a compact set in $(\mathbb{R}, \tau)$}{$(\mathbb{R}, \tau)$ is a first-countable space}{$(\mathbb{R}, \tau)$ is a second-countable space}
\ans{(A), (C) (MSQ)}

\textbf{Q.51 (MSQ)} Let $T, S : \mathbb{R}^4 \to \mathbb{R}^4$ be two non-zero, non-identity $\mathbb{R}$-linear transformations. Assume $T^2 = T$. Which of the following is/are TRUE?
\optslistA{$T$ is necessarily invertible}{$T$ and $S$ are similar if $S^2 = S$ and $\operatorname{Rank}(T) = \operatorname{Rank}(S)$}{$T$ and $S$ are similar if $S$ has only $0$ and $1$ as eigenvalues}{$T$ is necessarily diagonalizable}
\ans{(B), (D) (MSQ)}

\textbf{Q.52 (MSQ)} Let $p_1 < p_2$ be the two fixed points of the function $g(x) = e^x - 2$, where $x \in \mathbb{R}$. For $x_0 \in \mathbb{R}$, let the sequence $(x_n)_{n \ge 1}$ be generated by the fixed point iteration
\[
  x_n = g(x_{n-1}), \quad n \ge 1.
\]
Which one of the following is/are correct?
\optslistA{$(x_n)_{n \ge 0}$ converges to $p_1$ for any $x_0 \in (p_1, p_2)$}{$(x_n)_{n \ge 0}$ converges to $p_2$ for any $x_0 \in (p_1, p_2)$}{$(x_n)_{n \ge 0}$ converges to $p_2$ for any $x_0 > p_2$}{$(x_n)_{n \ge 0}$ converges to $p_1$ for any $x_0 < p_1$}
\ans{(A), (D) (MSQ)}

\textbf{Q.53 (MSQ)} Which of the following is/are eigenvalue(s) of the Sturm--Liouville problem
\[
  y'' + \lambda y = 0, \quad 0 \le x \le \pi, \quad y(0) = y'(0), \quad y(\pi) = y'(\pi)?
\]
\optsfourA{$\lambda = 1$}{$\lambda = 2$}{$\lambda = 3$}{$\lambda = 4$}
\ans{(A), (D) (MSQ)}

\textbf{Q.54 (MSQ)} Let $f : \mathbb{R}^2 \to \mathbb{R}$ be a function such that
\[
  f(x, y) = \begin{cases} \left(1 - \cos\dfrac{x^2}{y^2}\right)\sqrt{x^2 + y^2}, & \text{if } y \ne 0, \; x \in \mathbb{R}, \\ 0, & \text{otherwise.} \end{cases}
\]
Which of the following is/are correct?
\optslistA{$f$ is continuous at $(0, 0)$, but not differentiable at $(0, 0)$}{$f$ is differentiable at $(0, 0)$}{All the directional derivatives of $f$ at $(0, 0)$ exist and they are equal to zero}{Both the partial derivatives of $f$ at $(0, 0)$ exist and they are equal to zero}
\ans{(A), (D) (MSQ)}

\textbf{Q.55 (MSQ)} For an integer $n$, let $f_n(x) = x e^{-nx}$, where $x \in [0, 1]$. Let $S := \{f_n : n \ge 1\}$. Consider the metric space $(C([0, 1]), d)$, where
\[
  d(f, g) = \sup_{x \in [0, 1]} \{|f(x) - g(x)|\}, \quad f, g \in C([0, 1]).
\]
Which of the following statement(s) is/are true?
\optslistA{$S$ is an equi-continuous family of continuous functions}{$S$ is closed in $(C([0, 1]), d)$}{$S$ is bounded in $(C([0, 1]), d)$}{$S$ is compact in $(C([0, 1]), d)$}
\ans{(A), (C) (MSQ)}

\textbf{Q.56 (MSQ)} Let $T : \mathbb{R}^4 \to \mathbb{R}^4$ be an $\mathbb{R}$-linear transformation such that $1$ and $2$ are the only eigenvalues of $T$. Suppose the dimensions of $\operatorname{Kernel}(T - I_4)$ and $\operatorname{Range}(T - 2I_4)$ are $1$ and $2$, respectively. Which of the following is/are possible (upper triangular) Jordan canonical form(s) of $T$?
\optsfourA{$\begin{pmatrix} 1 & 0 & 0 & 0 \\ 0 & 2 & 0 & 0 \\ 0 & 0 & 2 & 1 \\ 0 & 0 & 0 & 2 \end{pmatrix}$}{$\begin{pmatrix} 1 & 0 & 0 & 0 \\ 0 & 2 & 1 & 0 \\ 0 & 0 & 2 & 1 \\ 0 & 0 & 0 & 2 \end{pmatrix}$}{$\begin{pmatrix} 1 & 1 & 0 & 0 \\ 0 & 1 & 0 & 0 \\ 0 & 0 & 2 & 1 \\ 0 & 0 & 0 & 2 \end{pmatrix}$}{$\begin{pmatrix} 1 & 1 & 0 & 0 \\ 0 & 1 & 0 & 0 \\ 0 & 0 & 2 & 0 \\ 0 & 0 & 0 & 2 \end{pmatrix}$}
\ans{(A), (D) (MSQ)}

\subsection*{Q.57 -- Q.65 Numerical Answer Type (NAT), carry TWO marks each (no negative marking).}

\textbf{Q.57} Let $L^2([-1, 1])$ denote the space of all real-valued Lebesgue square-integrable functions on $[-1, 1]$, with the usual norm $\|\cdot\|$. Let $\mathcal{P}_1$ be the subspace of $L^2([-1, 1])$ consisting of all the polynomials of degree at most $1$. Let $f \in L^2([-1, 1])$ be such that $\|f\|^2 = \frac{18}{5}$, $\int_{-1}^{1} f(x)\,dx = 2$, and $\int_{-1}^{1} x f(x)\,dx = 0$. Then
\[
  \inf_{g \in \mathcal{P}_1} \|f - g\|^2 = \underline{\hspace{2cm}} \quad \text{(round off to TWO decimal places)}.
\]
\ans{$1.59 \text{ to } 1.61$}

\textbf{Q.58} The maximum value of $f(x, y, z) = 10x + 6y - 8z$ subject to the constraints
\begin{align*}
  5x - 2y + 6z &\le 20, \\
  10x + 4y - 6z &\le 30, \\
  x, y, z &\ge 0,
\end{align*}
is equal to \underline{\hspace{2cm}} (round off to TWO decimal places).
\ans{$56.66 \text{ to } 56.68$}

\textbf{Q.59} Let $K \subseteq \mathbb{C}$ be the field extension of $\mathbb{Q}$ obtained by adjoining all the roots of the polynomial equation $(X^2 - 2)(X^2 - 3) = 0$. The number of distinct fields $F$ such that $\mathbb{Q} \subseteq F \subseteq K$ is equal to \underline{\hspace{2cm}} (answer in integer).
\ans{$5$}

\textbf{Q.60} Let $H$ be the subset of $S_3$ consisting of all $\sigma \in S_3$ such that
\[
  \operatorname{Trace}(A_1 A_2 A_3) = \operatorname{Trace}(A_{\sigma(1)} A_{\sigma(2)} A_{\sigma(3)}),
\]
for all $A_1, A_2, A_3 \in M_2(\mathbb{C})$. The number of elements in $H$ is equal to \underline{\hspace{2cm}} (answer in integer).
\ans{$3$}

\textbf{Q.61} Let $r : [0, 1] \to \mathbb{R}^2$ be a continuously differentiable path from $(0, 2)$ to $(3, 0)$ and let $\mathbf{F} : \mathbb{R}^2 \to \mathbb{R}^2$ be defined by $\mathbf{F}(x, y) = (1 - 2y, 1 - 2x)$. The line integral of $\mathbf{F}$ along $r$, $\int \mathbf{F} \cdot d\mathbf{r}$, is equal to \underline{\hspace{2cm}} (round off to TWO decimal places).
\ans{$0.99 \text{ to } 1.01$}

\textbf{Q.62} Let $u = u(x, t)$ be the solution of the initial value problem
\begin{align*}
  \frac{\partial^2 u}{\partial t^2} - \frac{\partial^2 u}{\partial x^2} &= 0, \quad x \in \mathbb{R}, \; t > 0, \\
  u(x, 0) &= 0, \quad x \in \mathbb{R}, \\
  \frac{\partial u}{\partial t}(x, 0) &= \begin{cases} x^4(1 - x)^4, & \text{if } 0 < x < 1, \\ 0, & \text{otherwise.} \end{cases}
\end{align*}
If $\alpha = \inf\{t > 0 : u(2, t) > 0\}$, then $\alpha$ is equal to \underline{\hspace{2cm}} (round off to TWO decimal places).
\ans{$0.99 \text{ to } 1.01$}

\textbf{Q.63} The boundary value problem
\begin{align*}
  x^2 y'' - 2x y' + 2y &= 0, \quad 1 \le x \le 2, \\
  y(1) - y'(1) &= 1, \\
  y(2) - k y'(2) &= 4,
\end{align*}
has infinitely many distinct solutions when $k$ is equal to \underline{\hspace{2cm}} (round off to TWO decimal places).
\ans{$1.99 \text{ to } 2.01$}

\textbf{Q.64} The global maximum of $f(x, y) = (x^2 + y^2) e^{2 - x - y}$ on $\{(x, y) \in \mathbb{R}^2 : x \ge 0, \, y \ge 0\}$ is equal to \underline{\hspace{2cm}} (round off to TWO decimal places).
\ans{$3.99 \text{ to } 4.01$}

\textbf{Q.65} Let $k \in \mathbb{R}$ and $D = \{(r, \theta) : 0 < r < 2, \, 0 < \theta < \pi\}$. Let $u(r, \theta)$ be the solution of the following boundary value problem
\begin{align*}
  \frac{\partial^2 u}{\partial r^2} + \frac{1}{r} \frac{\partial u}{\partial r} + \frac{1}{r^2} \frac{\partial^2 u}{\partial \theta^2} &= 0, \quad (r, \theta) \in D, \\
  u(r, 0) = u(r, \pi) &= 0, \quad 0 \le r \le 2, \\
  u(2, \theta) &= k \sin(2\theta), \quad 0 < \theta < \pi.
\end{align*}
If $u\!\left(1, \frac{\pi}{4}\right) = 2$, then the value of $k$ is equal to \underline{\hspace{2cm}} (round off to TWO decimal places).
\ans{$7.99 \text{ to } 8.01$}

\end{document}
